\documentclass[11pt]{article}
\usepackage[margin=1in]{geometry}
\usepackage[T1]{fontenc}
\usepackage{lmodern}
\usepackage[hidelinks]{hyperref}
\usepackage{booktabs}
\usepackage{microtype}

\title{Concrete Proof-Rule Cards: V3, V5, V6, and E1}
\author{Tactix Development Project}
\date{August 8, 2026}

\begin{document}
\maketitle

\section*{Executive summary}
The interactive proof-rule guide's Substitute, Sweep, Congruence, and BVE
entries have been replaced with concrete indexed transitions, ordered proof
deltas, reconstruction obligations, and checkable worked examples.  Their
authoritative entries in \path{docs/calculus.md} were tightened where the
existing prose did not name the emitted objects or their order.  This is a
documentation-only change; no solver behavior or performance changed.

\section{Findings and results}
\textbf{V3 Substitute.}  The rule now defines the SCC representative map, one
transient BIG-path bridge per rewritten polarity, the normalized image of each
touched clause, add-before-delete ordering, and the symmetric witness pair used
to recover each substituted variable.  The contradictory-SCC branch names both
RUP units and the final empty-clause chain.

\textbf{V5 Sweep.}  The environment now records how root-false literals are
stripped, and the two equivalence directions are named separately.  The guide
distinguishes main-kernel RUP replay from translation of a sub-solver derivation
and states the all-or-nothing validation rule.  A worked binary example gives
the exact hint order for both directions.

\textbf{V6 Congruence.}  The single ``equivalence ladders'' placeholder has
been split into the actual proof families: direct AND chains with BIG bridges,
the four ITE condition cases, XOR sign-pattern ladders, and degeneracy ladders.
The common order---transients, final conclusions, reverse transient
deletions---is explicit.  The examples show exact duplicate-AND chains and a
gate that becomes constant only after an earlier equivalence merge.

\textbf{E1 BVE.}  Positive and negative parents, indexed resolvents, root-unit
stripping, the fixed two-parent LRAT order, gate-reduced pair scope, and
both-sided reconstruction witnesses are now explicit.  The examples cover one
real resolvent and a tautological pair that emits no resolvent but still needs
witnesses.

\section{Evidence}
The revisions were checked against the corresponding implementations in
\path{src/rules/substitute}, \path{src/rules/sweep},
\path{src/rules/congruence}, and \path{src/rules/bve}.  The offline page parses
through Lynx and retains exactly 24 rule cards, 24 formal panels, 24 informal
panels, and 24 example panels.  Citation, section-reference, document-size,
comment-hygiene, rule-policy-boundary, and whitespace checks pass.  No solver
binary was changed, so no solver or performance measurement is attributed to
this increment.

\section{Risks and decisions}
The HTML remains a derived teaching view; \path{docs/calculus.md} remains
authoritative.  The expanded cards are necessarily denser, but the existing
tabs keep the formal detail separate from the plain-language explanation and
examples.  Three untouched cards still contain compact placeholders:
Vivify's ``reason-chain'', HyperBinaryRes's optional deletion, and Instantiate's
``reason-chain including C''.  They were left outside this four-rule scope
rather than revised without a dedicated semantic audit.

\section{Next actions}
Review the four cards in a graphical browser.  If the level of detail is right,
the remaining three shorthand cards can receive the same focused treatment.
The solver-development handoff remains unchanged: resume Kissat-default parity
with the \texttt{SP-COMPACT-1} row.

\end{document}
